\par\SourceLeaf{421}%
\noindent\hspace*{1.2em}Die Seite, nach welcher die Curve convex, oder die Zeichen in obigen Formeln werden also beim Fortgehen in der Curve so lange immer ungeändert bleiben als
\[
\frac{\delta V}{\delta\varrho}
\]
sein Zeichen nicht ändert. In so fern $V$ eine stetige Function ist, kann eine solche Änderung des Zeichens nur nach einem Durchgang durch den Werth $0$ erfolgen; dieser aber setzt nothwendig voraus, dass zugleich $p$ und $q=0$ werden. An einer solchen Stelle wird der Krümmungshalbmesser unendlich gross, oder die Krümmung verschwindet und in so fern dann, allgemein zu reden, hier
\[
-p\sin\varphi+q\cos\varphi
\]
sein Zeichen ändern wird, findet hier ein Wendungspunkt statt.

\begin{center}5.\end{center}
\noindent\hspace*{1.2em}Der Fall, wo die Natur der Curve dadurch ausgedrückt wird, dass $y$ einer gegebenen Function von $x$ gleich gesetzt wird, nemlich $y=X$, ist in dem vorigen enthalten, wenn man
\[
V=X-y
\]
setzt. Macht man
\[
dX=X'dx,\qquad dX'=X''dx,
\]
so wird
\[
\begin{gathered}
p=X',\qquad q=-1,\\
P=X'',\qquad Q=0,\qquad R=0,
\end{gathered}
\]
also
\[
\pm\frac1r=\frac{X''}{(1+X'X')^{\frac32}}.
\]
Da hier $q$ negativ ist, so gilt für wachsende $x$ das obere Zeichen. Man kann sich also hier kurz so ausdrücken, dass bei einem positiven $X''$ die Curve nach derselben Seite ihre Concavität hat, nach welcher die Axe der $y$ gegen die Axe der $x$ liegt; während bei negativem $X''$ die Curve nach dieser Seite zu convex ist.
